% !TEX root = ../../tree_notes.tex
% =====================================================================
%  4.1  Pointers vs Arrays                                  sec:c04-reps
%  Exports: def:c04-rep  def:c04-forms  thm:c04-implicit
%           lst:c04-reps  ex:c04-storage
%           prob:c04-nary-depth  prob:c04-validate-array  prob:c04-inform
%           lst:c04-nary-depth   lst:c04-validate-array   lst:c04-inform
% =====================================================================
\section{Pointers vs Arrays}\label{sec:c04-reps}

\begin{topicintro}
A tree is an abstract relation -- ``these nodes are the children of that node.'' To
compute on it you must commit that relation to memory, and there is more than one way
to do it. This section catalogues the \term{representations} a competitive programmer
actually meets -- \term{pointer nodes}, \term{adjacency lists}, \term{parent arrays}
and \term{implicit array indices} -- and, more importantly, \emph{when each wins}. The
same DFS (\chref{c03}) runs on top of every one; only the way you fetch a node's
children changes.
\textbf{Prerequisites.} \secref{c01-what-is-tree} (rooted vs unrooted, $\degop$), the
\code{TreeNode} model of \secref{c02-model}, the adjacency-build loop \coderef{c01-build}.
The two running examples $R$ (\figref{c04-running-r}) and $G$ (\figref{c04-running-g}).

\textbf{After this section you can}
\begin{itemize}[nosep]
  \item name the four storage models and fetch a node's children under each;      \hfill\lv{1}
  \item read an input format (edge list / \code{parent[]} / \code{leftChild[]}) and
        know which model it hands you;                                            \hfill\lv{2}
  \item convert a parent array into a children adjacency in one loop;             \hfill\lv{2}
  \item decide, from the operation mix, whether pointers or arrays are cheaper
        (cache, memory, upward navigation);                                       \hfill\lv{3}
  \item check whether an array-described graph is \emph{one} valid binary tree.    \hfill\lv{3}
\end{itemize}
\end{topicintro}

\begin{recognitionbox}[Pattern card -- how the tree is \emph{stored}]
\reckey{the problem's \emph{input} is a bunch of parallel arrays (\code{leftChild},
\code{rightChild}, \code{parent}), $n{-}1$ undirected \emph{edges}, or a pointer
\code{TreeNode}/\code{Node}; the difficulty is getting from a node to its children (or
parent) \emph{at all}, not what you do once you have them.}
\recuse{
  \item you are handed an \emph{edge list} or \code{parent[]} and must ``make a tree''
        before any recursion makes sense;
  \item the operation is ``go to the parent'' a lot -- pointers give $\bigO{1}$ only if
        you \emph{add} a \code{parent} pointer; arrays give it free;
  \item the tree is complete / heap-shaped and a flat array with $2i{+}1$ indexing is
        leaner than allocating nodes.
}
\recnotuse{
  \item the structure is given and the real work is the \emph{visit order} -- that is
        Chapter~3 (\chref{c03});
  \item you need subtree-as-contiguous-range times \code{tin}/\code{tout} -- Euler tour,
        \chref{c24};
  \item the array is a \emph{priority queue} with sift-up/down -- that is heaps,
        \chref{c07} (only the layout is teased here).
}
One-line test: \emph{is my bottleneck ``I cannot cheaply reach a node's children or
parent''? If yes, this is a representation choice -- this section.}
\end{recognitionbox}

\subsection{Intuition}

Forget the pictures for a moment and ask a machine question: node $u$ says ``I have
children'' -- \emph{where} are they written down? There are exactly four honest answers,
and every contest tree problem uses one of them.

\begin{figure}[H]
\centering
\begin{tikzpicture}[diagram, node distance=3mm,
   m/.style={boxnode, text width=39mm, minimum height=11mm, font=\scriptsize, align=left}]
  \node[m, fill=primaryblue!10, draw=primaryblue] (p) at (0,0)
    {\textbf{Pointer nodes}\\ \code{TreeNode*}/\code{Node*}\\ children are raw addresses\\
     \emph{up navigation: none}};
  \node[m, fill=accentorange!10, draw=accentorange, right=of p] (a)
    {\textbf{Adjacency list}\\ \code{vector<int> adj[N]}\\ children $=$ neighbours\\
     \emph{undirected; needs a root}};
  \node[m, fill=successgreen!10, draw=successgreen, right=of a] (par)
    {\textbf{Parent array}\\ \code{par[i] = p}\\ one entry per node\\
     \emph{up: $\bigO{1}$; down: rebuild}};
  \node[m, fill=softpurple!10, draw=softpurple, below=8mm of p] (im)
    {\textbf{Implicit index}\\ node $i$'s kids at $2i{+}1,2i{+}2$\\ zero pointers\\
     \emph{needs a \emph{complete} tree}};
  \draw[flow] (p) -- (a); \draw[flow] (a) -- (par);
\end{tikzpicture}
\caption{The four storage models. They describe the \emph{same} tree; they differ in
what a ``child'' physically is, whether you can step \emph{upward} in $O(1)$, and how
much memory the edges cost. Section~\ref{sec:c04-rooting} adds the fifth act: choosing
which node is the root once you have an adjacency list.}\label{fig:c04-models}
\end{figure}

The tension that generates every trade-off is simple. A \emph{pointer} tree knows its
children but not its parent, and every node is a separate heap allocation -- random-ish
addresses that scatter the data across RAM and defeat the CPU cache. An \emph{array}
tree keeps everything contiguous (great cache, tiny constants) but only expresses a
\emph{fixed arity} cleanly, and an implicit index $2i{+}1$ silently assumes the tree is
\emph{complete}. An \emph{adjacency list} is the general, contest-native compromise:
one contiguous vector per node, arbitrary arity, but \emph{undirected} -- which is why
rooting (\secref{c04-rooting}) is inseparable from it. Reading this chapter, hold one
question: \emph{which navigation (down? up? both?) does the problem need, and how many
edges can I afford to store?}

\subsection{Definitions and Notation}

\begin{definition}{Tree representation}{c04-rep}
A \term{representation} of a tree $T$ on $n$ nodes is an injective encoding of the
``is-a-child-of'' relation into a data structure supporting, at minimum, a
\texttt{children(v)} query returning $v$'s children in some order. A
representation is \term{rooted} if children are directed away from a distinguished root,
and \term{unrooted} if it stores the symmetric adjacency (neighbours) instead.
\end{definition}

\noindent\textbf{In plain words.} A representation is just \emph{where the edges live} so
that ``give me $v$'s children'' is a cheap, well-defined operation. All
four models below implement that one query; picking a model is picking the price of that
query and of the operations you can \emph{not} afford.

\begin{definition}{The four concrete models}{c04-forms}
For nodes labelled $0,\dots,n-1$:
\begin{itemize}[nosep]
  \item \textbf{Pointer}: \code{struct Node \{ int val; vector<Node*> kids; \};} -- a
        binary node instead stores \code{left},\ \code{right} (\code{nullptr} $=$ absent).
  \item \textbf{Adjacency list}: \code{vector<int> adj[n];} with each undirected edge
        $(u,v)$ stored \emph{both} ways: \code{adj[u].push\_back(v); adj[v].push\_back(u);}.
  \item \textbf{Parent array}: \code{int par[n];} with \code{par[root] = -1}; every other
        node stores its single parent.
  \item \textbf{Implicit index}: node $i$ has children at array slots $2i{+}1$ and
        $2i{+}2$ (and parent at $\lfloor (i-1)/2\rfloor$); the tree \emph{is} the array.
\end{itemize}
\end{definition}

\begin{notebox}[Adjacency \emph{list} vs \emph{matrix}]
A contest tree is \emph{sparse} ($n-1$ edges), so an adjacency \emph{matrix} ($O(n^2)$
memory) is never worth it; the list (or the parent array) is the default. Matrix form
only reappears for dense \emph{graph} adjacency in special chapters, not trees.
\end{notebox}

\subsection{Key Property: what implicit indexing costs}

\begin{theorem}{Implicit array indexing is exact only for complete trees}{c04-implicit}
Place a binary tree in an array $A$ with the root at index $0$ and node $i$'s
children at $2i+1$ (left) and $2i+2$ (right), for every non-leaf. Then this layout stores
the tree with no wasted slots \textbf{iff} the tree is \term{complete} (every level full
except possibly the last, which is filled left-to-right). A tree of height $h$ that is
\emph{not} complete can force $\size$ to occupy a slot $\ge 2^{h}-1$, so the array needs
up to $2^{h}$ entries to hold $h{+}1$ nodes.
\end{theorem}
\begin{proofbox}
By induction on the level. A node at index $i$ sits at depth $d=\lfloor\log_2(i+1)\rfloor$:
base $i=0$ has depth $0$; each step to $2i{+}1$ or $2i{+}2$ sends $i+1 \to 2(i+1)$ or
$2(i+1)+1$, doubling the ``offset $i{+}1$'', so depth rises by exactly $1$. Thus depth $d$
occupies indices $[\,2^{d}-1,\ 2^{d+1}-2\,]$, precisely the $2^{d}$ slots a full level
needs. A \emph{complete} tree fills each level left-to-right, so the $k$-th node (in
level order) is written to the $k$-th slot with zero gaps, and the count of slots equals
the node count. Conversely a right-leaning chain $0\to 2\to 5\to\dots$ (always the right
child) drives the index by $i\mapsto 2i+2$, i.e.\ slot $\ge 2^{d}$ at depth $d$, giving an
array of size $2^{\height{}}$ for only $\height{}+1$ nodes -- exponentially wasteful.
\end{proofbox}
\begin{insightbox}
This is the whole reason heaps (\chref{c07}) use the array layout and general trees do
not: a heap is \emph{complete by definition}, so $2i{+}1/2i{+}2$ is free navigation with
no pointers at all. The moment the tree can skew, that guarantee collapses -- which is
exactly what happens to running example $R$ (\figref{c04-running-r}), whose node $3$ has
no left child and therefore punches a \emph{hole} in the implicit array.
\end{insightbox}

\subsection{The Algorithm: fetching children under each model}

\begin{lstlisting}[caption={The same ``visit all children'' loop, four storages}, label={lst:c04-reps}]
#include <bits/stdc++.h>
using namespace std;

// (A) POINTER nodes -- binary, LeetCode's default. Up-navigation is IMPOSSIBLE
//     unless you pass the parent down or pre-store a parent pointer.
struct TreeNode { int val; TreeNode *left, *right;
    TreeNode(int x=0): val(x), left(nullptr), right(nullptr) {} };
void dfsPtr(TreeNode* v) {
    if (!v) return;
    // children are just the two fields; "none" is nullptr
    dfsPtr(v->left); dfsPtr(v->right);
}

// (B) ADJACENCY LIST -- general tree, contest default. UNDIRECTED, so you MUST
//     pass the parent to avoid bouncing back along the edge you came from.
vector<vector<int>> adj;              // adj[u] = neighbours of u
void dfsAdj(int u, int p) {
    for (int v : adj[u]) if (v != p)   // the parent guard: skip where we came from
        dfsAdj(v, u);
}

// (C) PARENT ARRAY par[i] = parent of i (par[root] = -1). Cheap to GO UP, but to
//     go DOWN you first rebuild children lists in one linear pass.
vector<int> par;                       // given
vector<vector<int>> buildKids(int n) {
    vector<vector<int>> kids(n);
    for (int i = 0; i < n; ++i)
        if (par[i] != -1) kids[par[i]].push_back(i);   // O(n), the key idiom
    return kids;
}

// (D) IMPLICIT INDEX -- the tree IS an array A; children of i are 2i+1, 2i+2.
//     Valid / compact only for a COMPLETE tree (thm:c04-implicit).
vector<int> A;                         // level-order packing
void dfsImpl(int i, int n) {
    if (i >= n) return;                // past the array = no node
    // visit A[i]; then
    dfsImpl(2*i + 1, n);               // left child slot
    dfsImpl(2*i + 2, n);               // right child slot
}
\end{lstlisting}
\begin{complexitybox}
All four walks are $\bigO{n}$ \emph{work} (each node once) plus the $\bigO{h}$ stack.
The representation changes only the \textbf{constants and side costs}: pointer $=$ one
heap allocation per node, scattered, no parent; adjacency $=$ contiguous per node but
$2(n{-}1)$ stored entries (each edge twice) and needs the parent guard; parent array $=$
$\bigO{1}$ up, $\bigO{n}$ to rebuild down; implicit index $=$ zero pointers, best cache,
but up to $O(2^{h})$ slots (thm:c04-implicit). Memory for the four: $O(n)$, $O(n)$,
$O(n)$, $O(2^{h})$ respectively.
\end{complexitybox}

\subsection{Worked Example}

\begin{example}{$R$ and $G$ in every model}{c04-storage}
\textbf{Binary $R$ (\figref{c04-running-r}).} As pointers, node $1$ has \code{left}
$=\&2$, \code{right} $=\&3$; node $3$'s \code{left} is \code{nullptr}. As an
implicit array with root slot $0$: $1$ at $0$, $2$ at $1$, $3$ at $2$, $4$ at $3$, $5$ at
$4$, and node $6$ (right child of $3$) at $2\cdot2{+}2 = 6$ -- so slot $5$ is a
\textbf{hole}, exactly as \thmref{c04-implicit} predicts ($R$ is not complete).

\medskip
\noindent\textbf{General $G$ (\figref{c04-running-g}).} Adjacency (undirected):
$\code{adj} = \{\{1,2,3\},\{0,4,5\},\{0\},\{0,6\},\{1\},\{1\},\{3\}\}$. Rooted at $0$
(\secref{c04-rooting}) the children are $0{:}\{1,2,3\},\ 1{:}\{4,5\},\ 3{:}\{6\}$, giving
parent array $\parent = [-1,0,0,0,1,1,3]$ and subtree sizes $\size=[7,3,1,2,1,1,1]$. Note
$\degop(1)=3$ in the \emph{undirected} graph equals $|\text{neighbours}|$, but its
\emph{child} count after rooting is only $2$ -- degree and child-count diverge the moment
a node has a parent, a distinction that trips up \probref{c04-validate-array}.
\end{example}
\begin{tracebox}[Parent array $\to$ children lists, one pass on $G$]
Start with \code{kids} $=$ seven empty vectors. Read $\parent=[-1,0,0,0,1,1,3]$:
$i{=}0$: $\parent[0]{=}-1$, skip (root). $i{=}1$: push $1$ under $0$. $i{=}2$: push $2$
under $0$. $i{=}3$: push $3$ under $0$. $i{=}4$: push $4$ under $1$. $i{=}5$: push $5$
under $1$. $i{=}6$: push $6$ under $3$. Result:
$\code{kids}=\{\{1,2,3\},\{4,5\},\{\},\{6\},\{\},\{\},\{\}\}$ -- $G$'s rooted
children, built in $n$ steps (\coderef{c04-reps}, block C). \checkmark\ \probref{c04-inform}.
\end{tracebox}

\subsection{Variations and Twists}

\begin{itemize}[nosep]
  \item \textbf{N-ary pointer}: \code{vector<Node*> children} instead of fixed
        \code{left}/\code{right} -- the model behind LC 559 (\probref{c04-nary-depth}). \hfill\lv{1}
  \item \textbf{Add a parent pointer} to a binary \code{TreeNode} to get $\bigO{1}$
        up-navigation (turn it into a quasi-linked structure) -- LC 2385 practice. \hfill\lv{3}
  \item \textbf{Edge list with a root given}: adjacency $+$ parent guard is the standard
        contest shape; the root is decided for you (\secref{c04-rooting}).          \hfill\lv{3}
  \item \textbf{Compressed adjacency} (\code{head[]/to[]/next[]} forward-star) to kill
        the per-node vector overhead when $n \sim 10^{6}$.                       \hfill\lv{5}
\end{itemize}

\subsection{Common Pitfalls}\label{sec:c04-reps-pitfalls}

\begin{warningbox}[Trap 1 -- forgetting the parent guard on an \emph{undirected} adjacency list]
\code{adj} stores each edge twice, so a bare DFS recurses $u \to v \to u \to v$ forever
(and overflows the stack). The guard \code{if (v == p) continue} is load-bearing.
\end{warningbox}
\begin{lstlisting}[caption={\textcolor{warningred}{WRONG}: bounces back along the edge}, language=C++, basicstyle=\ttfamily\small]
void dfs(int u) {
    for (int v : adj[u]) dfs(v);       // v can be u's parent -> infinite recursion
}
\end{lstlisting}
\begin{lstlisting}[caption={\textcolor{successgreen}{RIGHT}: carry the parent, skip it}, language=C++, basicstyle=\ttfamily\small]
void dfs(int u, int p) {
    for (int v : adj[u]) if (v != p) dfs(v, u);   // orient away from where we came
}
\end{lstlisting}

\begin{warningbox}[Trap 2 -- assuming implicit $2i{+}1$ indexing is compact for \emph{any} tree]
A skewed tree drives a node to slot $2^{d}$ at depth $d$, so an ``array of size $n$''
silently overflows or misreads. Implicit indexing is safe \emph{only} for complete trees
(\thmref{c04-implicit}); otherwise store real node indices in a vector, not heap slots.
\end{warningbox}
\begin{lstlisting}[caption={\textcolor{warningred}{WRONG}: right-skew blows the index up}, language=C++, basicstyle=\ttfamily\small]
A[2*i + 2] = x;        // depth d lands near slot 2^d -- array explodes on a chain
\end{lstlisting}
\begin{lstlisting}[caption={\textcolor{successgreen}{RIGHT}: use 2i indexing ONLY because a heap is complete}, language=C++, basicstyle=\ttfamily\small]
// guarantee: tree is complete (heap / full level-order) -> 2i+1, 2i+2 fit in n slots
\end{lstlisting}

\begin{warningbox}[Trap 3 -- conflating undirected $\degop$ with rooted child-count]
After rooting, a non-root node has $\text{children} = \degop - 1$ (the edge to its parent
is not a child edge). Counting children as $\degop$ mislabels leaves and breaks
leaf/root detection -- the exact bug in \probref{c04-validate-array} if you are careless.
\end{warningbox}
\begin{lstlisting}[caption={\textcolor{warningred}{WRONG}: degree is not child count}, language=C++, basicstyle=\ttfamily\small]
bool isLeaf = (adj[u].size() == 0);   // a rooted LEAF may still have 1 neighbour (its parent)
\end{lstlisting}
\begin{lstlisting}[caption={\textcolor{successgreen}{RIGHT}: count children after the guard}, language=C++, basicstyle=\ttfamily\small]
int kids = 0; for (int v : adj[u]) if (v != p) ++kids;
bool isLeaf = (kids == 0);
\end{lstlisting}

\subsection{Problem Lab}

Two model-reading exercises then a genuine validation problem. LC 559 is a $\lv{2}$
warm-up that is pure ``the children are a vector''; LC 1376 is a $\lv{3}$ core whose
\emph{entire} difficulty is turning a parent array into children lists; LC 1361 is a
$\lv{3}$ core that hands you parallel child arrays and asks whether they even describe
one tree -- a question only the storage model can answer.


\solvedproblem{c04-nary-depth}{Maximum Depth of N-ary Tree}{LeetCode 559}
  {https://leetcode.com/problems/maximum-depth-of-n-ary-tree/}{LC Easy L2}
  {N-ary pointer children vector, recursive depth}

\begin{problemstatement}[Maximum Depth of N-ary Tree]
Given the root of an \emph{N-ary} tree, return its depth -- the number of nodes on the
longest root-to-leaf path. Each node stores \code{int val} and \code{vector<Node*>
children} (an arbitrary, possibly empty, list). $0 \le n \le 10^4$, $0 \le \text{children}
\le 100$. An empty tree has depth $0$.
\end{problemstatement}

\solstep{Understanding the Problem}
This is Chapter~2's maximum-depth recursion (\secref{c02-bottomup}) with one
representation change: a node has a \emph{variable-length} list of children instead of a
fixed \code{left}/\code{right}. The whole exercise is reading model (A) of
\coderef{c04-reps} in its general form -- ``children'' is now a vector you loop over.

\solstep{Solution}
A leaf (empty \code{children}) has depth $1$; an internal node has depth $1 + \max$ over
its children's depths; \code{nullptr} gives $0$. One recursive pass, taking the max of the
child depths inside a \code{for}.

\solstep{Implementation}
\begin{lstlisting}[caption={Maximum Depth of an N-ary Tree (LeetCode 559)}, label={lst:c04-nary-depth}]
#include <bits/stdc++.h>
using namespace std;

class Node {
public:
    int val;
    vector<Node*> children;             // model (A), general arity
    Node() {}
    Node(int v) : val(v) {}
};

class Solution {
public:
    int maxDepth(Node* root) {
        if (!root) return 0;                        // empty tree has depth 0
        int best = 0;
        for (Node* c : root->children)              // loop replaces the 2 fixed calls
            best = max(best, maxDepth(c));          // deepest child so far
        return best + 1;                            // +1 for this node
    }
};
\end{lstlisting}

\solstep{Trace}
\begin{tracebox}[$R$ read as an N-ary tree: root $1$ has children $2,3$; $2$ has $4,5$; $3$ has $6$]
\code{maxDepth(4)}$=1$, \code{maxDepth(5)}$=1$ $\Rightarrow$ \code{maxDepth(2)}$=1{+}\max(1,1)=2$.
\code{maxDepth(6)}$=1$ $\Rightarrow$ \code{maxDepth(3)}$=2$. \code{maxDepth(1)}$=1{+}\max(2,2)=3$.
Depth $=3$ nodes ($1,3,6$), i.e.\ $\height(R){=}2$ edges $+1$. \checkmark\ \probref{c04-nary-depth}.
\end{tracebox}

\solstep{Complexity}
\begin{complexitybox}
Time $\bigO{n}$ (each node once), space $\bigO{h}$ recursion ($\bigO{n}$ worst case, a
chain). The iterative BFS-by-level alternative (\chref{c03}) uses $\bigO{w}$ width.
Verdict: \ac.
\end{complexitybox}

\begin{insightbox}[Key takeaway]
Nothing about the \emph{representation} (N-ary \code{children} vector) changes the
recursion -- only the loop that fetches children does. ``$\max$ over children, then $+1$''
is \secref{c02-bottomup} verbatim; the arity is a storage detail.
\end{insightbox}

% ---------------------------------------------------------------------------
\solvedproblem{c04-inform}{Time Needed to Inform All Employees}{LeetCode 1376}
  {https://leetcode.com/problems/time-needed-to-inform-all-employees/}{LC Medium L3}
  {parent-array $\to$ children adjacency, max-path DFS}

\begin{problemstatement}[Time Needed to Inform All Employees]
There are $n$ employees, a \emph{head} with id \code{headID}, and a \textbf{parent array}
\code{manager} where \code{manager[i]} is the direct manager of employee $i$
($\code{manager}[\code{headID}] = -1$). Employee $i$ needs \code{informTime[i]} minutes to
inform \emph{all} of his direct reports (each report starts only after $i$ finishes;
reports proceed in \textbf{parallel}). Return the minutes until every employee is
informed. $1 \le n \le 10^5$, $\sum \text{informTime} \le 2\cdot10^4$ fits comfortably.
\end{problemstatement}

\begin{attemptgate}[Attempt it yourself first -- \lv{3}, give yourself 20 minutes]
A fair first attempt: notice the input is a \emph{parent} array and figure out the right
way to walk it.
\peek{the tree is given bottom-up (\code{manager[i]} points \emph{up}); what must you
build before a top-down DFS is even possible?}
\nudge{once you have children lists, what single number at each node combines the
parallel informing? think max, not sum -- why?}
\fullhint{one linear pass turns \code{manager[]} into \code{children[]}; then one DFS
returns, per node, $\text{informTime}[\text{node}] + \max$ over children. The answer is
that value at the head. Full build below.}
\end{attemptgate}

\solstep{Understanding the Problem}
In tree language, $\code{manager}[\,]$ is a \term{parent array} (\coderef{c04-reps},
model C): each node names its parent, and the head is the root ($\parent=-1$). Informing
flows \emph{down} from root to leaves, so we must first convert the upward pointers into
downward \code{children} lists, then ask: over all root-to-leaf paths, which one takes
the longest to inform? Because parallel branches overlap, the finish time is the
\textbf{maximum} root-to-leaf path sum of $\code{informTime}$, not a total.

\solstep{Key Observations}
\begin{observation}
A parent array cannot be DFS'd top-down as given -- there is no way to enumerate a
manager's reports. The single idiom \code{children[manager[i]].push\_back(i)} (one pass,
$O(n)$) rebuilds a rooted children adjacency (\figref{c04-running-g}'s $\children$ lists).
\end{observation}
\begin{observation}
Reporting is a forest of chains but there is exactly one root (\code{headID}), so it is a
single tree. Informing is a longest-path problem: the last-informed employee is on the
root-to-leaf path maximising the sum of $\code{informTime}$.
\end{observation}
\begin{observation}
$\code{informTime}[\text{leaf}]$ contributes nothing to the answer (a leaf informs no
one), so leaves return $0$; only internal nodes add their broadcast time.
\end{observation}

\solstep{Think It Through With Me}
\begin{thinkbox}
\deadend{For each employee, follow \code{manager[]} up to the head summing times, then
take the max.}{This \emph{does} give the right max per employee, but naively costs
$O(n \cdot h)$ -- on a chain manager that is $O(n^2)$ and TLEs at $n=10^5$. Fixable by
memoising ``time-to-head'', but there is a cleaner $O(n)$ orientation we are missing.}
\deadend{Sum all $\code{informTime}$ as the total elapsed time.}{Dead wrong even at
$O(n)$: informing is \textbf{parallel}. If the head spends $t$ minutes then two reports
each broadcast to their teams \emph{simultaneously} -- those subtrees overlap in time.
The answer is the single slowest chain, a \emph{max}, never a sum.}
\aha{Orient the parent array into children lists once ($O(n)$), then one DFS returns
$\text{informTime}[u] + \max_{c}\text{dfs}(c)$; that value at \code{headID} is the answer.}
\end{thinkbox}

\solstep{The Approach Ladder}
\begin{center}\renewcommand{\arraystretch}{1.3}
\begin{tabular}{@{}c p{0.235\linewidth} p{0.28\linewidth} l c@{}}
\toprule
\textbf{\#} & \textbf{Approach} & \textbf{Idea} & \textbf{Time} & \textbf{Verdict} \\
\midrule
1 & \lv{2} sum all times & treat informing as serial & $O(n)$ & \textbf{WA} \\
2 & \lv{3} walk-up per employee & parent[] to root, max chain & $O(n\,h)$ & \tle \\
3 & \lv{3} orient $+$ max-path DFS & build children, one DFS, max & $O(n)$ & \ac \\
\bottomrule
\end{tabular}\end{center}
Rung 1 ignores parallelism (the model, not the maths, is wrong); rung 2 is right but
re-walks shared chains; rung 3 pays $O(n)$ once to flip the pointers and is the solution.

\solstep{Solution}
Build \code{children} by the model-C idiom. Then \code{dfs(u)} returns the slowest
informing time of $u$'s subtree: $0$ if $u$ is a leaf, else
$\code{informTime}[u] + \max_c \code{dfs}(c)$. Return $\code{dfs}(\code{headID})$.

\solstep{Proof of Correctness}
\begin{proofbox}[Induction on subtree height]
\emph{Claim:} $\code{dfs}(u)$ equals the minimum minutes from the moment $u$ starts
informing until the last employee in $\subtree(u)$ is informed. \emph{Base:} $u$ a leaf --
it informs no one, so $\code{dfs}(u)=0$. \emph{Step:} $u$ takes $\code{informTime}[u]$
minutes to reach all its direct reports; each report $c$ then independently informs
$\subtree(c)$, finishing $\code{dfs}(c)$ minutes after $c$ is notified. Since reports are
\textbf{parallel}, the subtree finishes when the \emph{slowest} report finishes:
$\code{informTime}[u] + \max_c \code{dfs}(c)$ -- exactly the code. By induction $\code{dfs}$
is correct at the root, and the root (\code{headID}) is where the cascade starts, so
$\code{dfs}(\code{headID})$ is the total elapsed time.
\end{proofbox}

\solstep{Implementation}
\begin{lstlisting}[caption={Time Needed to Inform All Employees (LeetCode 1376)}, label={lst:c04-inform}]
#include <bits/stdc++.h>
using namespace std;

class Solution {
public:
    int numOfMinutes(int n, int headID, vector<int>& manager, vector<int>& informTime) {
        // model C -> rooted children adjacency, one linear pass
        vector<vector<int>> children(n);
        for (int i = 0; i < n; ++i)
            if (manager[i] != -1) children[manager[i]].push_back(i);

        // slowest inform path in subtree(u): leaf -> 0
        function<int(int)> dfs = [&](int u) -> int {
            int best = 0;
            for (int c : children[u]) best = max(best, dfs(c));
            return best + informTime[u];          // + broadcast time of u
        };
        return dfs(headID);                        // headID is the root (manager == -1)
    }
};
\end{lstlisting}

\solstep{Trace}
\begin{tracebox}[$G$ re-read as the org chart, headID $=0$]
$\parent=[-1,0,0,0,1,1,3]$, $\code{informTime}=[1,2,3,4,0,0,0]$.
Children: $0\!\to\!\{1,2,3\}$, $1\!\to\!\{4,5\}$, $3\!\to\!\{6\}$.
Leaves $2,4,5,6$: $\code{dfs}(2)=3,\ \code{dfs}(4)=0,\ \code{dfs}(5)=0,\ \code{dfs}(6)=0$
(leaf $2$ still adds its own informTime).
$\code{dfs}(1)=2+\max(0,0)=2$; $\code{dfs}(3)=4+\max(0)=4$.
$\code{dfs}(0)=1+\max(\code{dfs}1{=}2,\code{dfs}2{=}3,\code{dfs}3{=}4)=1+4=\mathbf{5}$.
Answer $5$ -- the chain $0\!\to\!3\!\to\!6$ is slowest. \checkmark\ \probref{c04-inform}.
\end{tracebox}

\solstep{Complexity}
\begin{complexitybox}
Time $\bigO{n}$: one pass to build children, one DFS touching each node/edge once. Space
$\bigO{n}$ for the children lists plus $\bigO{h}$ recursion (worst-case $n$ on a chain --
see the hardening box). Fits $n=10^5$ with room.
\end{complexitybox}

\solstep{Verification and Edge Cases}
\begin{center}\renewcommand{\arraystretch}{1.2}
\begin{tabular}{@{}p{0.30\linewidth} p{0.62\linewidth}@{}}
\toprule
\textbf{Case} & \textbf{What the code does} \\
\midrule
$n=1$ (head only) & children empty, $\code{dfs(headID)}=0+\text{informTime}[\text{head}]$; a lone head informing nobody should give $\code{informTime}$ only if he had reports -- LC guarantees the answer here is the head's own time only when he reports; with no reports the loop is empty, returns $\code{informTime[head]}$ \\
star (head, all report to it) & $\max$ over many leaves each returning their own informTime; the answer is $\code{informTime[head]} + \max$ leaf \\
deep chain ($n$ in a line) & correct but $\bigO{h}$ recursion $=\bigO{n}$ -- RE risk, hardening below \\
manager[] lists head not at index 0 & root found purely by $\text{manager}[i]{=}-1$, index-independent \\
\bottomrule
\end{tabular}\end{center}

\solstep{Pitfalls and Variants}
\begin{warningbox}
(i) Trying to DFS \code{manager[]} directly -- it points \emph{up}, so you can never
enumerate reports; the children-list rebuild is mandatory. (ii) Summing instead of
max-ing (rung 1). (iii) Adding $\code{informTime}$ at a \emph{leaf}: a leaf informs
nobody, but the code still adds $\code{informTime[leaf]}$ -- that is \emph{correct} here
because LC sets $\code{informTime[leaf]}=0$; do not assume so in a variant with non-zero
leaf times. Variant: report each employee's finish time = the same DFS
threading an accumulated path-sum downward (a top-down parameter, \secref{c02-topdown}).
\end{warningbox}

\begin{insightbox}[Key takeaways -- three things to remember]
\begin{enumerate}[nosep, leftmargin=1.6em]
  \item A \textbf{parent array} is a top-down-blind representation: your first move is
        almost always \code{children[par[i]].push\_back(i)}, one $O(n)$ pass.
  \item \emph{Parallel} work on subtrees turns a sum into a \textbf{max} -- the depth of
        the process, not its total volume.
  \item The root is wherever the parent pointer is $-1$; never assume it is node~$0$.
\end{enumerate}
\end{insightbox}

\begin{hardeningbox}
\begin{itemize}[nosep, leftmargin=1.4em]
  \item \textbf{Chain $\Rightarrow$ recursion RE}: at $n=10^5$ a line-shaped org drives
        the \code{dfs} stack to $\bigO{n}$. Swap to an \emph{iterative} post-order (explicit
        stack, or process nodes in reverse BFS order once you have levels) -- the children
        array is unchanged, only the walk is.
  \item \textbf{Memory}: \code{vector<vector<int>> children} is $n$ separate small
        allocations; for $n \sim 10^6$ a forward-star \code{head/to/nxt} or a single flat
        array with an offset table beats it.
  \item \textbf{Overflow}: here $\sum\text{informTime}$ is bounded and fits \code{int};
        if a variant lets each of $10^5$ nodes add up to $10^9$, switch to \code{long long}.
\end{itemize}
\end{hardeningbox}

% ---------------------------------------------------------------------------
\solvedproblem{c04-validate-array}{Validate Binary Tree Nodes}{LeetCode 1361}
  {https://leetcode.com/problems/validate-binary-tree-nodes/}{LC Medium L3}
  {parallel child arrays, indegree + reachability, one-tree test}

\begin{problemstatement}[Validate Binary Tree Nodes]
You have $n$ nodes labelled $0\dots n-1$. Node $i$'s two children are
\code{leftChild[i]} and \code{rightChild[i]} (a value of $-1$ means no such child). Return
\code{true} \textbf{iff} the given arrays describe \emph{exactly one valid binary tree}.
$1 \le n \le 10^4$. Example: $n=4$, $\code{leftChild}=[1,-1,3,-1]$,
$\code{rightChild}=[2,-1,-1,-1]$ is \code{true} (root $0$ with left $1$, right $2$; $1$ has
left $3$).
\end{problemstatement}

\begin{attemptgate}[Attempt it yourself first -- \lv{3}, give yourself 25 minutes]
A fair first attempt: think about how many parents a single node may have.
\peek{in a \emph{tree}, how many nodes can point \emph{at} node $i$? What about the root?}
\peek{``one valid tree'' vs ``a forest of valid trees'' -- which extra check separates them?}
\nudge{compute the in-degree of every node from the two arrays; what does exactly one
in-degree-$0$ node with all others in-degree-$1$ fail to guarantee?}
\fullhint{Require every node in-degree $\le 1$, exactly one in-degree-$0$ root, and a
reachability DFS from that root visiting all $n$ nodes (which also kills cycles). Build
the DFS directly from \code{leftChild/rightChild} -- that \emph{is} the child adjacency.}
\end{attemptgate}

\solstep{Understanding the Problem}
The input \emph{is} a representation -- model (A) as two \textbf{parallel arrays}: instead
of pointers, node $i$'s children are stored by \emph{index} in
\code{leftChild[i]}/\code{rightChild[i]}. ``Is this one valid tree?'' is a structural
question about that array: (i) no node has two parents, (ii) there is a single root, and
(iii) everything is reachable from that root (no disconnected component, no cycle). Note
we are \emph{not} asked to build pointers -- the arrays already encode the same edges.

\solstep{Key Observations}
\begin{observation}
Each entry \code{leftChild[i]} $\neq -1$ declares an edge $i \to \code{leftChild[i]}$, so
the target's in-degree rises by $1$. A node with in-degree $>1$ has two parents -- not a
tree. A node with in-degree $0$ is a candidate root; a valid single tree has \textbf{exactly
one} such node.
\end{observation}
\begin{observation}
In-degree $\le 1$ everywhere and exactly one root are \emph{necessary but not sufficient}:
a two-node component that is a directed \emph{cycle} can leave an unrelated node as the
only in-degree-$0$ vertex, and a second valid-but-disconnected tree also passes those
two counts. Reachability is the missing axiom.
\end{observation}
\begin{observation}
A DFS (or BFS) from the unique root that never revisits a node and counts $n$ visited
nodes simultaneously proves connectivity, acyclicity, and the no-two-parents property --
so it can substitute for the explicit in-degree scan if you then also assert exactly one
root. We keep both checks: they are each one line and catch distinct bugs.
\end{observation}

\solstep{Think It Through With Me}
\begin{thinkbox}
\deadend{Check every node has in-degree $\le 1$ and stop.}
{$n=4$, $\code{leftChild}=[1,0,-1,-1]$, $\code{rightChild}=[-1,-1,3,-1]$: in-degrees are
$\{1,1,0,1\}$ -- all $\le 1$ -- yet node $0$ and node $1$ point at each other (a $2$-cycle),
and node $3$ hangs off $2$ which hangs from nowhere reachable from a root. Valid counts,
invalid tree: there is \textbf{no} in-degree-$0$ root that reaches them.}
\deadend{Count in-degree-$0$ nodes; accept iff there is exactly one.}{Add to the above a
genuinely disconnected second tree, e.g.\ edges $0\to1$ and $2\to3$ with $n=4$: there are
\textbf{two} in-degree-$0$ nodes ($0$ and $2$), so this is a \emph{forest}, not one tree.
Exactly-one-root is necessary but, as the cycle case shows, still not sufficient.}
\aha{One root \emph{and} full reachability: run a DFS from the unique in-degree-$0$ node,
reject any revisit (cycle / two-parent) and finally demand \code{visited}$=n$. That single
walk closes all three gaps at once.}
\end{thinkbox}

\solstep{The Approach Ladder}
\begin{center}\renewcommand{\arraystretch}{1.3}
\begin{tabular}{@{}c p{0.235\linewidth} p{0.28\linewidth} l c@{}}
\toprule
\textbf{\#} & \textbf{Approach} & \textbf{Idea} & \textbf{Time} & \textbf{Verdict} \\
\midrule
1 & \lv{2} in-degree $\le 1$ only & every node has $\le 1$ parent & $O(n)$ & \textbf{WA} \\
2 & \lv{2} exactly one in-degree-$0$ & a single root exists & $O(n)$ & \textbf{WA} \\
3 & \lv{3} root $+$ reachability DFS & one root, DFS visits all $n$, no revisit & $O(n)$ & \ac \\
\bottomrule
\end{tabular}\end{center}
Rungs 1 and 2 each check one structural axiom; the fourth (connectivity/acyclicity)
needs an actual traversal -- rung 3.

\solstep{Solution}
Compute in-degrees by scanning both child arrays. Reject if any in-degree $>1$. Find all
in-degree-$0$ nodes; reject unless there is exactly one (the root). DFS from the root over
\code{leftChild/rightChild}; reject if a node is entered twice (cycle); accept iff exactly
$n$ nodes were visited.

\solstep{Proof of Correctness}
\begin{proofbox}
A directed graph is a valid rooted binary tree on $n$ nodes iff (a) one node has in-degree
$0$ (the root), (b) every other node has in-degree $1$, and (c) all $n$ nodes are reachable
from the root. \emph{(a),(b):} the code rejects any in-degree $>1$ and any count of
in-degree-$0$ nodes $\neq 1$, so at most one root and at most one parent per node. \emph{(c):}
the DFS from the root visits the root-connected component; a \textbf{revisit} would mean a
node entered along two paths, i.e.\ a cycle or a second parent -- both rejected
immediately. If the walk ends having visited \code{seen}$=n$, every node is reachable from
the unique root and none is repeated, so the edges form a single acyclic connected
structure with correct in-degrees -- exactly a valid binary tree. Conversely any valid tree
satisfies (a)--(c), passes both guards and yields $\code{seen}=n$. Hence the procedure is
true on exactly the valid inputs.
\end{proofbox}

\solstep{Implementation}
\begin{lstlisting}[caption={Validate Binary Tree Nodes (LeetCode 1361)}, label={lst:c04-validate-array}]
#include <bits/stdc++.h>
using namespace std;

class Solution {
public:
    bool validateBinaryTreeNodes(int n, vector<int>& leftChild, vector<int>& rightChild) {
        vector<int> indeg(n, 0);
        for (int i = 0; i < n; ++i) {
            if (leftChild[i]  != -1) ++indeg[leftChild[i]];   // edge i -> left
            if (rightChild[i] != -1) ++indeg[rightChild[i]];  // edge i -> right
        }
        for (int d : indeg) if (d > 1) return false;          // some node has two parents

        int root = -1, roots = 0;
        for (int i = 0; i < n; ++i) if (indeg[i] == 0) { root = i; ++roots; }
        if (roots != 1) return false;                         // zero (cycle) or many (forest)

        vector<char> vis(n, 0);                               // reachability from the root
        stack<int> st; st.push(root); int seen = 0;
        while (!st.empty()) {
            int u = st.top(); st.pop();
            if (vis[u]) return false;                         // entered twice -> cycle
            vis[u] = 1; ++seen;
            if (leftChild[u]  != -1) st.push(leftChild[u]);
            if (rightChild[u] != -1) st.push(rightChild[u]);
        }
        return seen == n;                                      // one tree over all nodes
    }
};
\end{lstlisting}

\solstep{Trace}
\begin{tracebox}[$R$ re-coded as parallel arrays (the running binary tree)]
$n=6$, $\code{leftChild}=[1,3,-1,-1,-1,-1]$, $\code{rightChild}=[2,4,-1,-1,-1,5]$.
In-degrees: node $0{:}0$; $1{:}1$ (by $0$); $2{:}1$ (by $0$); $3{:}1$ (by $1$); $4{:}1$ (by $1$);
$5{:}1$ (by $2$). All $\le 1$; exactly one in-degree-$0$ (node $0$) $\Rightarrow$ root $=0$.
DFS from $0$ pushes children $1,2$, then $3,4$ and $5$ -- visits $0,1,2,3,4,5$ with no
revisit, $\code{seen}=6=n$ $\Rightarrow$ \textbf{true}. \checkmark\ \probref{c04-validate-array}.
\end{tracebox}

\solstep{Complexity}
\begin{complexitybox}
Time $\bigO{n}$: one scan to fill in-degrees, one scan for the root, one DFS touching each
edge once. Space $\bigO{n}$ (\code{indeg}, \code{vis}, explicit \code{stack}). Iterative
stack avoids any recursion-depth $\bigO{n}$ risk. Verdict: \ac.
\end{complexitybox}

\solstep{Verification and Edge Cases}
\begin{center}\renewcommand{\arraystretch}{1.2}
\begin{tabular}{@{}p{0.30\linewidth} p{0.62\linewidth}@{}}
\toprule
\textbf{Case} & \textbf{What the code does} \\
\midrule
$n=1$, empty child arrays & indeg$\{0\}$, one root, DFS sees $1$ $=n$ $\Rightarrow$ true \\
two parents & indeg of that node $=2$ $\Rightarrow$ immediate false \\
cycle $0\leftrightarrow1$ & some node has indeg $\ge1$ but no valid single root reaches all; revisit or seen$\neq n$ $\Rightarrow$ false \\
forest $0\to1,\ 2\to3$ & two in-degree-$0$ nodes $\Rightarrow$ roots$\neq1$ false \\
self-loop \code{leftChild[i]=i} & indeg[i] increments, and DFS would revisit $\Rightarrow$ false \\
\bottomrule
\end{tabular}\end{center}

\solstep{Pitfalls and Variants}
\begin{warningbox}
The classic wrong answer is stopping at ``all in-degrees $\le 1$ and one root'' (rung 2) --
it accepts a cycle-with-a-bystander-root and accepts nothing about connectivity. The
reachability DFS is not optional. Also, do not ``find the root by $\degop$'': root here is
the unique \textbf{in-degree-$0$} node, a directed notion -- a node with undirected degree
$1$ can be a leaf, not the root (cf.\ the degree-vs-children trap in
Trap~3 of \secref{c04-reps}). Variant: return the root index too (needed before a
downward traversal); the same in-degree scan finds it for free.
\end{warningbox}

\begin{insightbox}[Key takeaways -- three things to remember]
\begin{enumerate}[nosep, leftmargin=1.6em]
  \item ``Is this a tree?'' $=$ ``one root $+$ everyone has exactly one parent $+$
        all reachable'' -- three axioms, and a single DFS can enforce the last two.
  \item Parallel child arrays \emph{are} an adjacency representation; you never need to
        build \code{TreeNode}s to reason about the structure.
  \item Root $=$ in-degree $0$ (directed), \emph{not} a degree heuristic -- mixing
        undirected degree with directed parentage is the most common source of wrong answers here.
\end{enumerate}
\end{insightbox}

\begin{hardeningbox}
\begin{itemize}[nosep, leftmargin=1.4em]
  \item \textbf{Deep $+$ wide}: the explicit \code{stack} keeps you safe for a $10^4$-long
        chain where a recursive DFS would risk stack overflow; keep it iterative even when
        recursion is shorter to write.
  \item \textbf{Union-Find alternative}: the same three axioms can be checked by counting
        components ($=1$) and edges ($=n-1$) with a DSU -- see \chref{c28}; equivalent
        $O(n\,\alpha)$, useful if the input is a genuine undirected \emph{edge} list.
  \item \textbf{Two-parent subtlety}: an in-degree $>1$ is caught before the DFS, but a
        node that is its \emph{own} child also survives in-degree $\le1$ -- the DFS
        revisit check is what actually rejects it, so keep both.
\end{itemize}
\end{hardeningbox}

\subsection{Practice Problems}

\begin{stuckbox}
\begin{itemize}[nosep, leftmargin=1.4em]
  \item Given a \code{parent[]}/\code{manager[]} array -- have you built \code{children[]}
        \emph{before} attempting any top-down DFS? (Trap 1.)
  \item On an undirected adjacency list, does every DFS carry the parent guard?
  \item If the input is edge/child arrays, is your ``is it a tree?'' check enforcing
        \emph{reachability}, not just in-degree counts? (LC 1361 rung 2 vs rung 3.)
  \item Are you counting \emph{children} after rooting, or mistakenly using undirected
        $\degop$? (Trap 3.)
\end{itemize}
\end{stuckbox}

\begin{practiceset}
  \practice{Minimum Time to Collect All Apples in a Tree}{LeetCode 1443}
    {https://leetcode.com/problems/minimum-time-to-collect-all-apples-in-a-tree/}{LC Medium L3}
    {The input is $n{-}1$ \emph{undirected} edges: build \code{adj} with both directions,
    root at $0$ with a parent guard, and a subtree contributes $2$ per edge \emph{iff} it
    (or some descendant) holds an apple. Model (B) + one bottom-up DFS.}
  \practice{Amount of Time for Binary Tree to Be Infected}{LeetCode 2385}
    {https://leetcode.com/problems/amount-of-time-for-binary-tree-to-be-infected/}{LC Medium L3}
    {Infection spreads to parent \emph{and} children, but a \code{TreeNode} has no upward
    pointer -- first \textbf{augment} the pointer tree with a $\text{node}\to\text{parent}$
    map (or re-root pointers), then BFS from the start node; the answer is the farthest
    distance. The whole problem is a representation fix before any traversal.}
\end{practiceset}

\begin{recallbox}[10-second refresher -- 4.1 Pointers vs Arrays]
Four homes for the edges: \emph{pointer} (children $=$ addresses, no parent), \emph{adjacency
list} (undirected, parent guard required), \emph{parent array} (up $\bigO{1}$, rebuild
children to go down), \emph{implicit index} ($2i{+}1/2i{+}2$, only for complete trees).
Convert parent arrays with \code{children[par[i]].push\_back(i)}. Root $=$ in-degree $0$, not
degree.
\selftest{
  \item List the four storage models and the one query each makes \emph{cheap} and the one
        it makes \emph{impossible}.
  \item Write the parent-array $\to$ children-lists loop from memory (the \coderef{c04-reps}
        block C idiom).
  \item Why is implicit $2i{+}1$ indexing compact only for a \emph{complete} tree? Give the
        slot a right-chain at depth $d$ reaches.
  \item State the three axioms that make an array-described graph ``one valid tree'', and
        name the check that enforces all three at once.
  \item On $G$, build $\children$ and $\parent$ from \figref{c04-running-g} without looking.
}
\end{recallbox}
